\documentclass[12pt]{article}
\usepackage[utf8]{inputenc}
\usepackage[T1]{fontenc}
\usepackage{hyperref}
\title{Degano Language Resources for K-12 Curriculum Development}
\author{Compiled by Claude Sonnet 5.5 (claude-sonnet-5-5) for Andr\'as Kornai}
\date{}
\begin{document}
\maketitle

\section{Introduction}
ISO 639-3 code wsv corresponds, in the sources I found, to the language called Wotapuri-Katarqalai (also Wot\={a}p\={u}r\={\i}-Kat\={a}qal\={a}\={\i}, Wotapuri, Dardu, Katarkalai). I found no search result using the name ``Degano''. The identification rests on the ISO code and the Glottocode given in the request. The language is a Dardic or Nuristani-related Indo-European language of the Katar valley, Nuristan, Afghanistan. One source reports only three elderly male speakers in September 2023 (one in Katar-qala, two in Quro), while another gives an earlier extinction date. Its status is therefore uncertain and at best critically endangered.

\section{Primary Educational Resources}
I found no primers, textbooks, school programmes or curricula for this language.

\section{Linguistic and Reference Materials}
General reference entries exist, but no grammar or dictionary turned up in my searches.
\begin{itemize}
\item Wikipedia article: \url{https://en.wikipedia.org/wiki/Wotapuri-Katarqalai_language}
\item Endangered Languages Project entry: \url{https://endangeredlanguages.com/lang/3319}
\end{itemize}

\section{Religious and Community Resources}
I found none.

\section{Digital and Audio Resources}
I found no audio, video or corpus resources.

\section{Curriculum Development Context}
Very little is documented. A revitalization curriculum would have to start from fieldwork with the few remaining speakers and from archival wordlists and notes from the 1940s to 1960s, which I did not locate online. Any programme should be built with the Katar valley community.
\end{document}